% Reusable ICOS margin specimens for the Harbor Book.
%
% Ground truth for meanings: the local International Code of Signals corpus
% (Pub. 102, US ed. 1969, rev. 2003) in the ICOS skill data/signals.json.
% The five fixed colors below follow the existing repository's real-flag SVG
% renderer (website-v2/src/components/ui/SignalFlags.tsx).  These are flag
% faces, not Port Daddy semantic status colors.
%
% A registered two-letter group is drawn as one signal: its two flags touch.
% Separate signals in a longer hoist belong in \icosHoist and receive a
% tackline.  This prevents visual adjacency from inventing sentence grammar.

% The Book preamble already loads TikZ, xcolor, and the PGF libraries. This
% fragment is intentionally preamble-neutral so the Book remains the proof
% environment rather than a standalone font substitute.

\definecolor{icosblue}{HTML}{003399}
\definecolor{icosred}{HTML}{CC0000}
\definecolor{icosyellow}{HTML}{FFD100}
\definecolor{icosblack}{HTML}{000000}
\definecolor{icosink}{HTML}{18212B}
\definecolor{icosrule}{HTML}{8996A3}

\newlength{\icosflagwidth}
\newlength{\icosflagheight}
\setlength{\icosflagwidth}{0.27in}
\setlength{\icosflagheight}{0.18in}

\newcommand{\icosFlagFrame}{%
  \draw[icosrule, line width=.18pt] (0,0) rectangle (\icosflagwidth,\icosflagheight);%
}

\newcommand{\icosFaceA}{% Alpha: white hoist, blue fly, swallowtail notch.
  \pgfmathsetlengthmacro{\anotch}{0.22*\fw}
  \path[fill=white]
    (0,0) -- (\fw,0) -- ({\fw-\anotch},\fh*.5) -- (\fw,\fh) -- (0,\fh) -- cycle;
  \path[fill=icosblue]
    (\fw*.5,0) -- (\fw,0) -- ({\fw-\anotch},\fh*.5) -- (\fw,\fh) -- (\fw*.5,\fh) -- cycle;}
\newcommand{\icosFaceC}{% Charlie: blue-white-red-white-blue stripes.
  \path[fill=icosblue] (0,0) rectangle (\fw,\fh*.2);
  \path[fill=white] (0,\fh*.2) rectangle (\fw,\fh*.4);
  \path[fill=icosred] (0,\fh*.4) rectangle (\fw,\fh*.6);
  \path[fill=white] (0,\fh*.6) rectangle (\fw,\fh*.8);
  \path[fill=icosblue] (0,\fh*.8) rectangle (\fw,\fh);}
\newcommand{\icosFaceH}{% Hotel: vertically divided white and red.
  \path[fill=white] (0,0) rectangle (\fw*.5,\fh);
  \path[fill=icosred] (\fw*.5,0) rectangle (\fw,\fh);}
\newcommand{\icosFaceK}{% Kilo: vertically divided yellow and blue.
  \path[fill=icosyellow] (0,0) rectangle (\fw*.5,\fh);
  \path[fill=icosblue] (\fw*.5,0) rectangle (\fw,\fh);}
\newcommand{\icosFaceL}{% Lima: yellow and black quarters.
  \path[fill=icosyellow] (0,\fh*.5) rectangle (\fw*.5,\fh);
  \path[fill=icosblack] (\fw*.5,\fh*.5) rectangle (\fw,\fh);
  \path[fill=icosblack] (0,0) rectangle (\fw*.5,\fh*.5);
  \path[fill=icosyellow] (\fw*.5,0) rectangle (\fw,\fh*.5);}
\newcommand{\icosFaceQ}{% Quebec: plain yellow field.
  \path[fill=icosyellow] (0,0) rectangle (\fw,\fh);}
\newcommand{\icosFaceS}{% Sierra: white field with blue square.
  \path[fill=white] (0,0) rectangle (\fw,\fh);
  \path[fill=icosblue] (\fw*.22,\fh*.16) rectangle (\fw*.78,\fh*.84);}
\newcommand{\icosFaceX}{% X-ray: white field with blue cross.
  \path[fill=white] (0,0) rectangle (\fw,\fh);
  \path[fill=icosblue] (\fw*.42,0) rectangle (\fw*.58,\fh);
  \path[fill=icosblue] (0,\fh*.38) rectangle (\fw,\fh*.62);}
\newcommand{\icosFaceY}{% Yankee: yellow and red diagonal stripes.
  \path[fill=icosyellow] (0,0) rectangle (\fw,\fh);
  \begin{scope}
    \clip (0,0) rectangle (\fw,\fh);
    \foreach \x in {-12,-4,4,12,20,28,36} {
      \path[fill=icosred] (\x,0) -- (\x+5,0) -- (\x+\fh+5,\fh) -- (\x+\fh,\fh) -- cycle;
    }
  \end{scope}}
\newcommand{\icosFaceZ}{% Zulu: four triangles, yellow-blue-red-black.
  \path[fill=icosyellow] (0,\fh) -- (\fw*.5,\fh*.5) -- (\fw,\fh) -- cycle;
  \path[fill=icosblue] (\fw,\fh) -- (\fw*.5,\fh*.5) -- (\fw,0) -- cycle;
  \path[fill=icosred] (\fw,0) -- (\fw*.5,\fh*.5) -- (0,0) -- cycle;
  \path[fill=icosblack] (0,0) -- (\fw*.5,\fh*.5) -- (0,\fh) -- cycle;}

\makeatletter
\newcommand{\icosOutlineA}{%
  \pgfmathsetlengthmacro{\anotch}{0.22*\fw}
  \draw[icosrule, line width=.18pt]
    (0,0) -- (\fw,0) -- ({\fw-\anotch},\fh*.5) -- (\fw,\fh) -- (0,\fh) -- cycle;
}
\newcommand{\icosFlag}[1]{%
  \begingroup
    % The Book preamble wraps PGF's picture boxing to center full-width
    % artwork. Restore its saved raw boxer for these inline flag faces so each
    % face keeps its native width inside the margin specimen.
    \ifcsname pd@typesetpicturebox\endcsname
      \let\pgfsys@typesetpicturebox\pd@typesetpicturebox
    \fi
  \begin{tikzpicture}[x=1pt,y=1pt,baseline=-.34ex]
    \pgfmathsetlengthmacro{\fw}{\icosflagwidth}
    \pgfmathsetlengthmacro{\fh}{\icosflagheight}
    \csname icosFace#1\endcsname
    \ifcsname icosOutline#1\endcsname
      \csname icosOutline#1\endcsname
    \else
      \icosFlagFrame
    \fi
  \end{tikzpicture}%
  \endgroup
}
\makeatother

\newcommand{\icosRegisteredGroup}[1]{%
  \ifcsname icosGroup#1\endcsname
    \csname icosGroup#1\endcsname
  \else
    \icosFlag{#1}%
  \fi
}
\newcommand{\icosGroupAS}{\icosFlag{A}\icosFlag{S}}
\newcommand{\icosGroupCS}{\icosFlag{C}\icosFlag{S}}
\newcommand{\icosGroupZL}{\icosFlag{Z}\icosFlag{L}}

\newcommand{\icosTackline}{\hspace{.07in}\rule[-.02in]{.11in}{.35pt}\hspace{.07in}}

% Separate signals use this macro; each comma-separated argument is a complete
% signal and the tackline is the visual boundary.  Do not use it to invent a
% two-letter meaning absent from the corpus.
\newcommand{\icosHoist}[2]{\icosRegisteredGroup{#1}\icosTackline\icosRegisteredGroup{#2}}

% Margin content only: the Book's pdmarginexhibit title/caption slots are
% intentionally left empty in the integration patch. The flag, one exact
% translation, and one concrete Book connection are the complete specimen.
  \newcommand{\icosCompactSpecimen}[3]{%
    \begingroup
    % The Book's enclosing \marginfont supplies the native margin face and size.
    % Keep this macro content-only so it can grow vertically instead of shrinking.
    \raggedright
  \centering\textbf{#1}\enspace\icosRegisteredGroup{#1}\par
  \vspace{.045in}
  \raggedright
  \textit{#2}\par
  \vspace{.04in}
  #3\par
  \endgroup
}

